\documentclass[11pt]{article}

\usepackage[T1]{fontenc}
\usepackage[utf8]{inputenc}
\usepackage[margin=1in]{geometry}
\usepackage{lmodern}
\usepackage{microtype}
\usepackage{amsmath,amssymb,mathtools,amsthm}
\usepackage{aliascnt}
\usepackage{booktabs,tabularx,array,longtable,multirow}
\usepackage[section]{placeins}
\usepackage{float}
\usepackage{graphicx}
\usepackage{xcolor}
\usepackage{enumitem}
\usepackage{needspace}
\usepackage{url}
\usepackage[hidelinks]{hyperref}
\usepackage[nameinlink,noabbrev]{cleveref}
\usepackage[backend=biber,style=numeric,sorting=none,maxbibnames=99]{biblatex}
\addbibresource{references.bib}

\allowdisplaybreaks
\setlength{\emergencystretch}{2em}
\raggedbottom
\setlist{leftmargin=*,topsep=3pt,itemsep=1.5pt,parsep=0pt,partopsep=0pt}

\definecolor{deepblue}{RGB}{23,69,124}
\hypersetup{
  colorlinks=true,
  linkcolor=deepblue,
  citecolor=deepblue,
  urlcolor=deepblue,
  pdftitle={The Alignment Trap: Structural Limits of SNOVA-Style Matrix Signatures},
  pdfauthor={Justin Thaler, Kai Zhang, Scott Kominers, Quang Dao},
  pdfsubject={Alignment rigidity and Hermitian feature-capacity bounds for matrix-based MQ signatures},
  pdfkeywords={SNOVA, key-randomness alignment, double centralizer, normalizer, Hermitian quotient, multivariate signatures}
}

\newtheorem{theorem}{Theorem}[section]
\newaliascnt{lemma}{theorem}
\newtheorem{lemma}[lemma]{Lemma}
\aliascntresetthe{lemma}
\newaliascnt{proposition}{theorem}
\newtheorem{proposition}[proposition]{Proposition}
\aliascntresetthe{proposition}
\newaliascnt{corollary}{theorem}
\newtheorem{corollary}[corollary]{Corollary}
\aliascntresetthe{corollary}
\newaliascnt{definition}{theorem}
\newtheorem{definition}[definition]{Definition}
\aliascntresetthe{definition}
\newaliascnt{remark}{theorem}
\newtheorem{remark}[remark]{Remark}
\aliascntresetthe{remark}

\newcommand{\F}{\mathbb F}
\newcommand{\Mat}{\operatorname{Mat}}
\newcommand{\Sym}{\operatorname{Sym}}
\newcommand{\rank}{\operatorname{rank}}
\newcommand{\Span}{\operatorname{span}}
\newcommand{\im}{\operatorname{im}}
\newcommand{\transpose}{\mathsf T}
\newcommand{\SNOVA}{\textsf{SNOVA}}

\title{The Alignment Trap:\\Structural Limits of SNOVA-Style Matrix Signatures}
\author{%
  Justin Thaler\thanks{a16z crypto research and Georgetown University} \and
  Kai Zhang\thanks{MIT} \and
  Scott Kominers\thanks{Harvard University and a16z crypto research} \and
  Quang Dao\thanks{Carnegie Mellon University and LayerZero Labs}%
}
\date{29 July 2026}

\begin{document}
\maketitle

\begin{abstract}
Matrix-based MQ signatures can compress public keys by choosing secret block
coefficients that align with public left and right multiplier algebras.  We
show that full-capacity alignment is rigid.  If the multiplier algebras are
degree-$d$ fields and a $d$-dimensional linear secret support commutes with
them on the two sides, then the support and the multiplier fields are forced
to be one adjoint pair.  Allowing secret coefficients to relabel multipliers
by field automorphisms does not escape: every nonzero linear support contained
in a field normalizer lies in one Frobenius-semilinear coset.  Even linearity
is not essential: a sufficiently large arbitrary two-sided normalizer support
forces the same adjoint pair.  Replacing the field by a genuinely
noncommutative irreducible algebra only shrinks the dimension of its commuting
multiplier algebra.

For odd characteristic, the adjoint pair induces a canonical involution on
the $d^2$ ordered left/right labels.  Symmetric base forms factor through a
Hermitian quotient of dimension $d(d+1)/2$, while skew-symmetric forms factor
through the complementary anti-Hermitian quotient of dimension $d(d-1)/2$.
A transpose-stable $w$-dimensional base-form family therefore cannot realize
$wd^2$ independent ordered features.  Achieving the raw count requires a
maximally non-self-adjoint family whose symmetric and skew projections are
both injective and whose two polynomial image channels are full and disjoint.

These results give a precise trilemma for SNOVA-style designs: accept the
Hermitian capacity loss, restore a full anti-Hermitian/exterior channel, or
abandon the alignment-based congruence-reuse mechanism.  The theorem does not
claim that every construction in the latter two branches is insecure.  It
shows that none is a local repair preserving the present alignment philosophy
and its ordered-label security count.  As an application, the symmetric
$q=19$ SNOVA Version~2.3 forms collapse $16$ labels to $10$ when $d=4$ and
$4$ labels to $3$ when $d=2$; the resulting concrete all-parameter forgery is
developed separately in the companion cryptanalysis.
\end{abstract}

\section{Introduction}\label{sec:intro}

SNOVA-style matrix signatures obtain compact public descriptions by reusing a
small family of congruence-transformed quadratic forms.  Their public maps are
expanded by multiplying matrix-valued variables on the left and right by
members of a small algebra.  A security analysis may then count the ordered
left/right labels as distinct quadratic directions.  The same algebraic
alignment that makes the public-key derivation work, however, constrains which
label spaces can genuinely remain independent.

This paper isolates that architectural constraint.  It deliberately does not
repeat the concrete forgery analysis of SNOVA Version~2.3; that analysis is
given in the companion paper~\cite{ThalerEtAl2026SymmetryQuotient}.  Here the
question is broader and structural:

\begin{quote}
How much left/right feature capacity can a matrix-based MQ design retain while
its secret coefficients align with the public multiplier algebras strongly
enough to reuse one congruence-transformed family of forms?
\end{quote}

Our first answer is a capacity theorem.  Let $\mathcal T$ be an irreducible
secret algebra acting on $\F_q^d$, and let $D$ be its centralizer.  The finite
double-centralizer theorem forces
\[
  D\cong\F_{q^e},\qquad
  \mathcal T\cong\Mat_{d/e}(\F_{q^e})
\]
for a divisor $e\mid d$.  Every pointwise aligned multiplier algebra lies in
$D$ or $D^\transpose$, so it has dimension at most $e$.  The field case
$e=d$ is therefore the unique extremum retaining $d$ multiplier dimensions on
each side.  Moving to a genuinely noncommutative secret algebra enlarges the
secret algebra only by shrinking the public multiplier capacity.

At full multiplier dimension the rigidity is stronger.  Degree-$d$ field
algebras in $\Mat_d(\F_q)$ are self-centralizing.  Hence a full
$d$-dimensional pointwise aligned secret support is exactly the right
multiplier field and the transpose of the left multiplier field.  This does
not assume that the support is an algebra, is irreducible, or is closed under
multiplication.  Automorphism relabeling also fails to create a new degree of
freedom: a linear support whose nonzero elements normalize a degree-$d$ field
is contained in one Frobenius-semilinear coset.  We further give a cardinality
criterion showing that sufficiently large nonlinear supports are equally
rigid.

Our second answer concerns quadratic features.  Once alignment forces the
left and right fields to be an adjoint pair, scalar transposition defines a
canonical involution on their ordered tensor labels.  Symmetric base forms see
only the Hermitian quotient; skew-symmetric forms see only the anti-Hermitian
quotient.  This yields the exact capacities $d(d+1)/2$ and $d(d-1)/2$.
For a family $W$ of base matrices, the total rank is controlled by the ranks
of the symmetric and skew projections of $W$.  A transpose-stable family
always loses capacity.  Recovering all ordered labels requires
$W\cap W^\transpose=0$, full symmetric and skew projections, and disjoint
polynomial image channels.  Thus ``make the forms slightly nonsymmetric'' is
not a local repair.

The resulting conclusion is a trilemma rather than a universal impossibility
claim.  A design must accept the Hermitian cap, deliberately restore the full
anti-Hermitian channel, or change the alignment and key-derivation mechanism.
The second and third branches may contain secure constructions, but they are
architecturally different from the symmetric aligned design whose security is
based on counting ordered labels.

\paragraph{Contributions.}
\begin{enumerate}
\item We classify the multiplier capacity of irreducible aligned secret
      algebras by the finite double-centralizer theorem.
\item We prove rigidity for pointwise alignment, semilinear normalizer
      alignment, and sufficiently large nonlinear supports.
\item We derive the Hermitian and anti-Hermitian quotient dimensions for an
      arbitrary adjoint pair of multiplier fields.
\item We quantify the nonsymmetry necessary to recover the raw ordered-label
      count and formulate the resulting no-local-repair trilemma.
\item We specialize the theorem to the symmetric public forms of SNOVA
      Version~2.3 and connect the structural loss to the companion concrete
      cryptanalysis.
\end{enumerate}

\paragraph{Scope.}
The alignment theorems hold over arbitrary finite fields; the decomposition
into symmetric and anti-Hermitian channels uses odd characteristic.  We do not
assert that every maximally non-self-adjoint or non-aligned descendant is
broken.  We prove that such a descendant has left the design regime in which
full alignment, transpose-stable forms, and a $d^2$ ordered-label count coexist.

\section{Model and notation}\label{sec:model}
Let $V=\F_q^d$ and identify $\operatorname{End}_{\F_q}(V)$ with
$\Mat_d(\F_q)$.  A \emph{degree-$d$ field algebra} is an embedded copy
$A\cong\F_{q^d}$ in $\Mat_d(\F_q)$.  For a subalgebra $B$, write
\[
 \mathcal C(B)=\{X:XB=BX\}
\]
for its centralizer and $N(B)$ for the normalizer of $B^\times$ in
$\operatorname{GL}_d(\F_q)$.

A right multiplier algebra $A_R$ and a left multiplier algebra $A_L$ act on a
base matrix $P$ through ordered labels $\xi\otimes\eta\in A_L\otimes A_R$:
\[
  \Phi_P(\xi\otimes\eta)(u)
    =u^\transpose\Sigma_\xi P\Sigma_\eta u.
\]
The block embeddings $\Sigma_\xi$ do not affect the local $d\times d$
algebraic arguments, so we suppress their ambient identity factors when
studying alignment.

Pointwise alignment means that each secret block coefficient $t$ satisfies
\[
 t^\transpose a_L=a_Lt^\transpose,
 \qquad
 a_Rt=ta_R
\]
for all public multipliers.  Semilinear alignment relaxes commutation to
normalization, allowing conjugation by $t$ to relabel field elements by a
Frobenius automorphism.  These are exactly the identities needed to move
public multipliers past secret coefficients while keeping one
congruence-transformed base-form family reusable.

\section{Alignment capacity and rigidity}\label{sec:alignment-rigidity}
\label{sec:alignment-forces-adjoint}

The public-key derivation uses two properties of the alignment algebra at
once.  Secret block coefficients commute with the public right multipliers,
and their transposes commute with the public left multipliers.  These are not
basis conventions: with full-dimensional support, they determine the two
multiplier algebras.

The first theorem removes a possible escape before specializing to a field.
It is a finite double-centralizer statement: a noncommutative irreducible
secret algebra can be large only by making its commuting multiplier algebra
small.

\begin{theorem}[Alignment capacity and the field extremum]
\label{thm:alignment-capacity}
Let $\mathcal T\subseteq\Mat_d(\F_q)$ be a unital algebra acting
irreducibly on $V=\F_q^d$, and put $D=\mathcal C(\mathcal T)$.  Then there
is a divisor $e\mid d$ such that
\begin{equation}
 D\cong\F_{q^e},
 \qquad
 \mathcal T=\mathcal C(D)\cong
       \Mat_{d/e}(\F_{q^e}),
 \qquad
 \dim_{\F_q}\mathcal T=\frac{d^2}{e}.          \label{eq:double-centralizer-shape}
\end{equation}
If public right and left multiplier algebras satisfy
\[
 A_R\subseteq\mathcal C(\mathcal T),
 \qquad
 A_L\subseteq\mathcal C(\mathcal T^\transpose),
\]
then
\begin{equation}
 \dim_{\F_q}A_R\le e,
 \qquad
 \dim_{\F_q}A_L\le e.                         \label{eq:alignment-capacity}
\end{equation}
In particular, retaining a $d$-dimensional multiplier algebra on either side
forces $e=d$, $\dim_{\F_q}\mathcal T=d$, and
$\mathcal T=D\cong\F_{q^d}$.  Every genuinely noncommutative irreducible
choice has $e<d$ and therefore strictly fewer than $d^2$ ordered multiplier
labels before any symmetry quotient is taken.
\end{theorem}

\begin{proof}
By Schur's lemma, $D=\operatorname{End}_{\mathcal T}(V)$ is a finite
division algebra.  Wedderburn's little theorem makes it a field
$\F_{q^e}$.  The faithful $D$-action makes $V$ a $D$-vector space, so
$e\mid d$ and $r:=\dim_DV=d/e$.  Because the assumed irreducibility says
that $V$ is a simple $\mathcal T$-module, the Jacobson density theorem
identifies the image of $\mathcal T$ with $\operatorname{End}_D(V)$.
All spaces are finite, so density is equality~\cite{Jacobson1989BasicAlgebraII}.
Hence
$\mathcal T=\mathcal C(D)\cong\Mat_r(\F_{q^e})$ and has
$\F_q$-dimension $er^2=d^2/e$.

The two alignment inclusions place $A_R$ inside $D$ and $A_L$ inside
$D^\transpose=\mathcal C(\mathcal T^\transpose)$, proving
\eqref{eq:alignment-capacity}.  Since $e\le d$, a multiplier algebra of
$\F_q$-dimension $d$ forces $e=d$ and $r=1$.  Then
$\mathcal T=D$ is itself a degree-$d$ field.  If $\mathcal T$ is
noncommutative, $r>1$, so $e<d$ and the product of the left and right label
dimensions is at most $e^2<d^2$.
\end{proof}

For $d=4$, the alternatives are especially stark: the field extremum has a
four-dimensional centralizer and sixteen ordered labels; an irreducible
$\Mat_2(\F_{q^2})$ secret algebra has only a two-dimensional centralizer
and four ordered labels; and the full matrix algebra has only scalar
multipliers.  The irreducibility hypothesis belongs only to this capacity
classification.  It is not a loophole at full multiplier dimension:
\Cref{prop:pointwise-alignment-rigidity} below traps every pointwise aligned
secret coefficient in the degree-$d$ multiplier field, whether the generated
algebra is irreducible or reducible.

\begin{lemma}[A full-degree field is self-centralizing]
\label{lem:field-centralizer}
Let $B\subseteq\Mat_d(\F_q)$ be a field of $\F_q$-dimension $d$.  Then
\[
  \mathcal C(B)=B.
\]
\end{lemma}

\begin{proof}
View $V=\F_q^d$ as a $B$-module.  The embedding is faithful and
$\dim_{\F_q}B=\dim_{\F_q}V=d$, so $V$ is one-dimensional over $B$
(and hence the $B$-action is automatically irreducible).  An
$\F_q$-linear endomorphism commutes with $B$ exactly when it is
$B$-linear.  Every $B$-linear endomorphism of the
one-dimensional $B$-space $V$ is multiplication by an element of $B$.
\end{proof}

The field centralizer is also the building block for more permissive
alignment in which secret mixing relabels public multipliers by a field
automorphism rather than commuting with them pointwise.

\begin{lemma}[Normalizer and uniqueness of the aligned field]
\label{lem:field-normalizer}
Let $d>1$, let $B\subseteq\Mat_d(\F_q)$ be as in
\Cref{lem:field-centralizer}, and let $\varphi:x\mapsto x^q$ denote the
$\F_q$-linear Frobenius map after identifying $\F_q^d$ with $B$.  Then
\begin{equation}
 N(B)=B^\times\rtimes\langle\varphi\rangle.          \label{eq:singer-normalizer}
\end{equation}
Moreover, if $C\subseteq\Mat_d(\F_q)$ is another field of
$\F_q$-dimension $d$ and $C^\times\subseteq N(B)$, then $C=B$.
\end{lemma}

\begin{proof}
Conjugation by an element of $N(B)$ induces an $\F_q$-automorphism of $B$,
so it induces $x\mapsto x^{q^k}$ for some $k$.  Composing with
$\varphi^{-k}$ leaves an element that centralizes $B$, hence lies in
$B^\times$ by \Cref{lem:field-centralizer}.  This proves
\eqref{eq:singer-normalizer}.

Choose a generator $c$ of the cyclic group $C^\times$, so
$\operatorname{ord}(c)=q^d-1$.  Write $c=b\varphi^k$ using
\eqref{eq:singer-normalizer}.  If $k\ne0$, put
$g=\gcd(d,k)$ and $e=d/g$.  Then $(b\varphi^k)^e$ is a norm lying in
$\F_{q^g}^\times$, and therefore
\[
 \operatorname{ord}(c)\mid e(q^g-1)
 \le d(q^{d/2}-1)<q^d-1.
\]
The final strict inequality follows from $q^{d/2}+1>d$, which holds for
$q\ge2$ and $d>1$.
This contradicts $\operatorname{ord}(c)=q^d-1$.  Thus $k=0$, so
$c\in B^\times$ and $C=\F_q[c]=B$.
\end{proof}

A more permissive design might choose secret coefficients from a linear
space that normalizes the multiplier field, so that moving a multiplier past
a secret coefficient applies a Frobenius automorphism rather than fixing the
multiplier pointwise.  Even this freedom is rigid at the SNOVA dimensions.

\begin{lemma}[Linear normalizer support is one semilinear coset]
\label{lem:linear-normalizer-support}
Let $A\subseteq\Mat_d(\F_q)$ be a degree-$d$ field and let
$\varphi:x\mapsto x^q$ be Frobenius in an $A$-coordinate model.  If
$0\ne U\subseteq\Mat_d(\F_q)$ is an $\F_q$-linear subspace satisfying
\begin{equation}
 U\setminus\{0\}\subseteq N(A),                    \label{eq:linear-support-normalizes}
\end{equation}
then, for every $0\ne u_0\in U$,
\begin{equation}
 Uu_0^{-1}\subseteq A,
 \qquad
 U\subseteq Au_0=A\varphi^k\text{ for one }0\le k<d,
 \qquad
 \dim_{\F_q}U\le d.                               \label{eq:semilinear-coset-support}
\end{equation}
If $\dim U=d$, then $U=A\varphi^k$.  If additionally $1\in U$, then
$U\subseteq A$, with equality when $\dim U=d$.

Now let $U$ have dimension $d$.  If
$U\setminus\{0\}$ normalizes a right multiplier field $A_R$ and
$U^\transpose\setminus\{0\}$ normalizes a left multiplier field $A_L$,
then
\begin{equation}
  A_L=A_R^\transpose.                              \label{eq:linear-support-adjoint-pair}
\end{equation}
Thus full-dimensional automorphism-relabeling support still determines the
same adjoint multiplier pair; it can only twist the secret support by one
Frobenius-semilinear coset.  No relation between $q$ and $d$ is required.
\end{lemma}

\begin{proof}
The standard linearized-polynomial representation gives the direct-sum
decomposition
\begin{equation}
 \operatorname{End}_{\F_q}(A)
   =\bigoplus_{k=0}^{d-1}A\varphi^k.                \label{eq:linearized-direct-sum}
\end{equation}
Indeed, the maps $x\mapsto\sum_{k=0}^{d-1}a_kx^{q^k}$ are linearly
independent: a nonzero such polynomial has ordinary degree at most
$q^{d-1}<q^d$ and therefore cannot vanish on all $q^d$ elements of $A$.
Their total $\F_q$-dimension is $d^2$, so they span every $\F_q$-linear
endomorphism.  Together with \eqref{eq:singer-normalizer}, this says that
every nonzero normalizer element lies in exactly one summand
$A\varphi^k$.

Fix $0\ne u_0\in U$ and put $V=Uu_0^{-1}$.  Then $V$ is an
$\F_q$-linear subspace, $1\in V$, and
$V\setminus\{0\}\subseteq N(A)$.  Suppose that
$v=a\varphi^k\in V$ with $k\ne0$.  By the direct sum
\eqref{eq:linearized-direct-sum}, $1+v$ is nonzero and has nonzero
components in two distinct summands.  Hence $1+v\notin N(A)$, contradicting
$1+v\in V$.  Thus every element of $V$ lies in $A$, proving
$Uu_0^{-1}\subseteq A$.  Writing $u_0=a_0\varphi^k$ now gives
$U\subseteq A\varphi^k$ and all remaining one-sided conclusions.

For the two-sided statement, the first part and $\dim U=d$ give
$Uu_0^{-1}=A_R$.  Since $u_0$ normalizes $A_R$, this also gives
$u_0^{-1}U=A_R$.  Applied to $U^\transpose$, the same argument gives
$A_L=U^\transpose u_0^{-\transpose}$.  Therefore
\[
 A_L=U^\transpose u_0^{-\transpose}
     =(u_0^{-1}U)^\transpose=A_R^\transpose,
\]
which proves \eqref{eq:linear-support-adjoint-pair}.
\end{proof}

\begin{lemma}[Large nonlinear normalizer support is rigid]
\label{lem:large-normalizer-intersection}
Let $A,C\subseteq\Mat_d(\F_q)$ be degree-$d$ field algebras.  If
$A\ne C$, then
\begin{equation}
 |N(A)\cap N(C)|
 \le d^2\bigl(q^{d-1}-1\bigr).                    \label{eq:normalizer-intersection-bound}
\end{equation}
Consequently, let $\mathcal S\subseteq\Mat_d(\F_q)$ be an arbitrary
not-necessarily-linear support such that every nonzero $t\in\mathcal S$
normalizes a right multiplier field $A_R$ and $t^\transpose$ normalizes a
left multiplier field $A_L$.  If
\begin{equation}
 |\mathcal S\setminus\{0\}|>
 d^2\bigl(q^{d-1}-1\bigr),                         \label{eq:large-support-threshold}
\end{equation}
then $A_L=A_R^\transpose$.  In particular, a support of cardinality $q^d$
containing zero forces the adjoint pair whenever $q\ge d^2$.
\end{lemma}

\begin{proof}
Write
\[
 N(A)=\bigcup_{k=0}^{d-1}(A\varphi^k\setminus\{0\}),
 \qquad
 N(C)=\bigcup_{j=0}^{d-1}(C\psi^j\setminus\{0\})
\]
using \Cref{lem:field-normalizer}.  The displayed components are
$d$-dimensional $\F_q$-linear subspaces after adjoining zero.  If
$A\varphi^k=C\psi^j$ for some pair $(k,j)$, choose a nonzero element $u_0$
in that common space.  Right-normalizing the common space by $u_0^{-1}$
gives both $A$ and $C$, so $A=C$.  Hence, when $A\ne C$, every pairwise
intersection has dimension at most $d-1$ and contains at most
$q^{d-1}-1$ nonzero elements.  A union bound over the $d^2$ pairs proves
\eqref{eq:normalizer-intersection-bound}.

Put $C=A_L^\transpose$.  If $t^\transpose$ normalizes $A_L$, transposing
the conjugation identity shows that $t$ normalizes $C$.  Thus
$\mathcal S\setminus\{0\}\subseteq N(A_R)\cap N(C)$.  The strict inequality
\eqref{eq:large-support-threshold} forces $A_R=C=A_L^\transpose$.  Finally,
when $|\mathcal S|=q^d$ and $0\in\mathcal S$, the required strict inequality
follows from
\[
 q^d-1-d^2(q^{d-1}-1)
 =q^{d-1}(q-d^2)+(d^2-1)>0
\]
for $q\ge d^2$.
\end{proof}

\begin{proposition}[Pointwise full-degree alignment is rigid]
\label{prop:pointwise-alignment-rigidity}
Let $A_L,A_R\subseteq\Mat_d(\F_q)$ be degree-$d$ field algebras and let
$U\subseteq\Mat_d(\F_q)$ be an $\F_q$-linear secret-coefficient support.
If
\begin{equation}
 t^\transpose a_L=a_Lt^\transpose,
 \qquad
 a_Rt=ta_R                                      \label{eq:alignment-commutation}
\end{equation}
for every $t\in U$, $a_L\in A_L$, and $a_R\in A_R$, then
\begin{equation}
 U\subseteq A_R\cap A_L^\transpose,
 \qquad
 \dim_{\F_q}U\le d.                              \label{eq:pointwise-support-intersection}
\end{equation}
If $\dim_{\F_q}U=d$, then
\begin{equation}
  U=A_R=A_L^\transpose.                            \label{eq:pointwise-adjoint-pair}
\end{equation}
This conclusion needs neither irreducibility of the algebra generated by
$U$ nor closure of $U$ under multiplication.
\end{proposition}

\begin{proof}
By \Cref{lem:field-centralizer}, the degree-$d$ field $A_R$ is
self-centralizing.  The right commutation in
\eqref{eq:alignment-commutation} therefore puts every $t\in U$ in $A_R$.
The left commutation puts $t^\transpose$ in
$\mathcal C(A_L)=A_L$, hence $t\in A_L^\transpose$.  This proves
\eqref{eq:pointwise-support-intersection}.  If $\dim U=d$, equality of
dimensions forces \eqref{eq:pointwise-adjoint-pair}.
\end{proof}

The alignment condition is visible on one central summand.  Let $t$ be a
secret block coefficient, let $a_L$ and $a_R$ be public left and right
multipliers, and substitute $X=TU$.  One obtains terms of the form
\[
 U_i^\transpose t^\transpose a_L F a_R sU_j.
\]
To reuse one congruence-transformed public form while keeping the public
multipliers outside it, the derivation moves $a_L$ past $t^\transpose$ and
$a_R$ past $s$.  The official construction achieves this by drawing all of
these coefficients from one symmetric commutative field
\cite{SNOVA23}.  The general consequence is stronger.

\begin{theorem}[Full alignment determines the multiplier pair]
\label{thm:alignment-multiplier-pair}
Let $A_L,A_R\subseteq\Mat_d(\F_q)$ be degree-$d$ field algebras.
The multiplier algebras are necessarily an adjoint pair under any of the
following alignment hypotheses:
\begin{enumerate}[label=(\roman*)]
\item a $d$-dimensional $\F_q$-linear secret support $U$ satisfies
      pointwise alignment \eqref{eq:alignment-commutation};
\item a $d$-dimensional $\F_q$-linear support $U$ has every nonzero
      element normalizing $A_R$ and every nonzero element of $U^\transpose$
      normalizing $A_L$; or
\item an arbitrary support $\mathcal S$ satisfies the same two-sided
      normalizer condition and
      $|\mathcal S\setminus\{0\}|>d^2(q^{d-1}-1)$.
\end{enumerate}
More precisely,
\begin{equation}
  A_L=B^\transpose,
  \qquad
  A_R=B                                             \label{eq:adjoint-pair}
\end{equation}
for the degree-$d$ field $B=A_R$.  In case (i),
$U=B=A_L^\transpose$.  In case (ii),
$U=B\varphi^k$ for one Frobenius power $\varphi^k$.  Case (iii) requires
neither linearity nor algebra closure of the support.  No irreducibility or
$q$-versus-$d$ hypothesis is needed.
\end{theorem}

\begin{proof}
Case (i) is \Cref{prop:pointwise-alignment-rigidity}, case (ii) is
\Cref{lem:linear-normalizer-support}, and case (iii) is
\Cref{lem:large-normalizer-intersection}.  Each gives
$A_L=A_R^\transpose$; putting $B=A_R$ gives \eqref{eq:adjoint-pair}.
\end{proof}

Thus choosing different left and right algebras does not make them
independent: full alignment forces the adjoint pair $B^\transpose$ and $B$.
The official same-algebra construction is the special case in which this pair
coincides.

\begin{corollary}[Same-algebra alignment forces transpose stability]
\label{cor:normalizer-alignment}
If the secret coefficients and both multiplier families use one degree-$d$
field algebra $A$, then full alignment forces
\begin{equation}
 A^\transpose=A.                                    \label{eq:transpose-stable-A}
\end{equation}
If pointwise public multipliers span $A$ but one insists on
$A^\transpose\ne A$, every alignable secret coefficient lies in the proper
subspace $A\cap A^\transpose$.
\end{corollary}

\begin{proof}
The first claim is \eqref{eq:adjoint-pair} with $A_L=A_R=B=A$.  For the
second, an alignable $t\in A$ has
$t^\transpose\in\mathcal C(A)=A$, so
$t\in A\cap A^\transpose$.
\end{proof}

When \eqref{eq:transpose-stable-A} holds, transpose induces
\[
 \sigma:A\longrightarrow A,
 \qquad \sigma(a)=a^\transpose.
\]
Because $A$ is commutative, this anti-automorphism is an automorphism.

\begin{corollary}[Classification of the same-algebra adjoint]
\label{cor:alignment-adjoint-classification}
The map $\sigma$ is an $\F_q$-automorphism of
$A\cong\F_{q^d}$ satisfying $\sigma^2=1$.  Consequently
\[
 \sigma=\operatorname{id},
 \quad\text{or, when $d$ is even,}\quad
 \sigma(x)=x^{q^{d/2}}.
\]
In particular, when $d$ is odd, full same-algebra alignment forces every
element of $A$ to be symmetric.  Changing the generator or basis cannot
produce any other adjoint.
\end{corollary}

\begin{proof}
For $a,b\in A$,
$\sigma(ab)=b^\transpose a^\transpose=\sigma(a)\sigma(b)$, and transpose is
$\F_q$-linear and involutory.  The $\F_q$-automorphism group of
$\F_{q^d}$ is cyclic of order $d$, generated by Frobenius.  It has no
nontrivial element of order two when $d$ is odd and has the unique such
element $x\mapsto x^{q^{d/2}}$ when $d$ is even.
\end{proof}

\section{The Hermitian quotient forced by alignment}
\label{sec:alignment-hermitian}

Assume now that $q$ is odd.  Under full alignment, write
$A_L=B^\transpose$ and $A_R=B$ as in \eqref{eq:adjoint-pair}.  Define an
involution on the ordered left/right label space by
\begin{equation}
 \tau(\xi\otimes\eta)
   =\eta^\transpose\otimes\xi^\transpose,
 \qquad
 \xi\in A_L,\ \eta\in A_R.                         \label{eq:twisted-label-involution}
\end{equation}
The Hermitian and anti-Hermitian label quotients are
\begin{align}
 \operatorname{Her}^2(A_L,A_R)
   &= (A_L\otimes_{\F_q}A_R)/
      \Span\{z-\tau z:z\in A_L\otimes A_R\},       \label{eq:hermitian-square}\\
 \operatorname{AHer}^2(A_L,A_R)
   &= (A_L\otimes_{\F_q}A_R)/
      \Span\{z+\tau z:z\in A_L\otimes A_R\}.       \label{eq:antihermitian-square}
\end{align}
For a base matrix $P$, put
\[
 \Phi_P(\xi\otimes\eta)(u)
   =u^\transpose\Sigma_\xi P\Sigma_\eta u.
\]

\begin{theorem}[Alignment--Hermitian factorization]
\label{thm:alignment-hermitian-factorization}
For every base matrix $P$,
\begin{equation}
  \Phi_P\circ\tau=\Phi_{P^\transpose}.              \label{eq:transpose-intertwining}
\end{equation}
Consequently:
\begin{enumerate}[label=(\roman*)]
\item if $P^\transpose=P$, then $\Phi_P$ factors through
      $\operatorname{Her}^2(A_L,A_R)$;
\item if $P^\transpose=-P$, then $\Phi_P$ factors through
      $\operatorname{AHer}^2(A_L,A_R)$;
\item the quotient dimensions are
\begin{equation}
 \dim_{\F_q}\operatorname{Her}^2(A_L,A_R)
      =\frac{d(d+1)}2,
 \qquad
 \dim_{\F_q}\operatorname{AHer}^2(A_L,A_R)
      =\frac{d(d-1)}2.                              \label{eq:hermitian-dimensions}
\end{equation}
\end{enumerate}
Thus every fully aligned symmetric redesign has the same universal
$d(d+1)/2$ cap, even if it uses distinct left and right algebras; additional
dependencies can only lower the rank.
\end{theorem}

\begin{proof}
The block embedding is transpose-compatible:
$\Sigma_{a^\transpose}^\transpose=\Sigma_a$.  Therefore, for a pure
tensor, transposing the scalar gives
\begin{align*}
 \Phi_P(\tau(\xi\otimes\eta))(u)
 &=u^\transpose\Sigma_{\eta^\transpose}P
                  \Sigma_{\xi^\transpose}u\\
 &=u^\transpose\Sigma_\xi P^\transpose\Sigma_\eta u
 =\Phi_{P^\transpose}(\xi\otimes\eta)(u),
\end{align*}
which proves \eqref{eq:transpose-intertwining}.  The two factorization
statements follow.

Choose a basis $b_1,\ldots,b_d$ of $B$ and the corresponding basis
$b_1^\transpose,\ldots,b_d^\transpose$ of $B^\transpose$.  In the basis
$b_i^\transpose\otimes b_j$, the involution $\tau$ is ordinary swap:
\[
 b_i^\transpose\otimes b_j
 \longmapsto b_j^\transpose\otimes b_i.
\]
Because $q$ is odd, the tensor product is the direct sum of the $+1$ and
$-1$ eigenspaces.  They have dimensions $(d^2+d)/2$ and $(d^2-d)/2$,
which are the dimensions of the two quotients.
\end{proof}

\begin{corollary}[Noncommutative alignment only lowers symmetric capacity]
\label{cor:noncommutative-capacity}
Let $\mathcal T$ be any irreducible secret algebra from
\Cref{thm:alignment-capacity}, with centralizer
$D\cong\F_{q^e}$.  Under pointwise alignment using the full centralizer pair
$D^\transpose,D$, every symmetric base form has ordered-label rank at most
\begin{equation}
  \frac{e(e+1)}2.                                  \label{eq:general-centralizer-hermitian-cap}
\end{equation}
The raw label space has dimension only $e^2$.  Hence the field extremum
$e=d$ simultaneously maximizes the aligned multiplier capacity and maximizes
the symmetric quotient; every genuinely noncommutative irreducible secret
algebra has a strictly smaller cap.
\end{corollary}

\begin{proof}
Choose a basis $b_1,\ldots,b_e$ of $D$ and the transpose basis of
$D^\transpose$.  On the $e^2$ labels
$b_i^\transpose\otimes b_j$, scalar transposition again acts by ordinary
swap.  A symmetric form therefore identifies $(i,j)$ with $(j,i)$ and leaves
at most $e(e+1)/2$ directions.  The strict comparison follows from $e<d$ in
the noncommutative case of \Cref{thm:alignment-capacity}.
\end{proof}

\begin{table}[H]
\centering
\caption{The alignment-capacity alternatives for the public block sizes.
The final column is the symmetric-form cap when the full centralizer pair is
used.  Noncommutativity increases secret-algebra dimension only by shrinking
the public multiplier labels.}
\label{tab:alignment-capacity-examples}
\small
\renewcommand{\arraystretch}{1.12}
\begin{tabular}{@{}c l r r r r@{}}
\toprule
block $d$ & irreducible secret algebra $\mathcal T$ & $e$ & $\dim\mathcal T$ & raw labels & symmetric cap\\
\midrule
$2$ & $\F_{q^2}$ & $2$ & $2$ & $4$ & $3$\\
$2$ & $\Mat_2(\F_q)$ & $1$ & $4$ & $1$ & $1$\\
$4$ & $\F_{q^4}$ & $4$ & $4$ & $16$ & $10$\\
$4$ & $\Mat_2(\F_{q^2})$ & $2$ & $8$ & $4$ & $3$\\
$4$ & $\Mat_4(\F_q)$ & $1$ & $16$ & $1$ & $1$\\
\bottomrule
\end{tabular}
\end{table}

In the official same-algebra case $A_L=A_R=A$, the adjoint $\sigma$ from
\Cref{cor:alignment-adjoint-classification} rewrites
\eqref{eq:twisted-label-involution} as
\[
 \tau(\xi\otimes\eta)=\sigma(\eta)\otimes\sigma(\xi).
\]
When $\sigma=\operatorname{id}$, the two quotients are the ordinary symmetric
and alternating squares.  The theorem therefore rules out three superficial
repairs at once.  A nonsymmetric generator does not help if same-algebra
alignment is retained; distinct left and right algebras are forced to be an
adjoint pair and give the same quotient; and increasing $q$ does not change
the dimensions.

\section{How much nonsymmetry is required to escape}
\label{sec:maximally-nonselfadjoint}

The preceding theorem does not claim that every arbitrary nonsymmetric base
family is already broken.  It instead quantifies how radical such a repair
must be.  Let $W$ be a $w$-dimensional space of base matrices and define
\[
 W_+=\Span\{P+P^\transpose:P\in W\},
 \qquad
 W_-=\Span\{P-P^\transpose:P\in W\},
 \qquad s_\pm=\dim W_\pm.
\]
Let $\mathcal I_+$ and $\mathcal I_-$ be the polynomial image spaces
generated by $W_+$ and $W_-$, respectively, over all aligned ordered labels.

\begin{theorem}[Hermitian-channel rank bound]
\label{thm:mixed-base-rank-bound}
For the complete feature map generated by $W$ and the aligned ordered
label space $A_L\otimes A_R$,
\begin{equation}
 \rank\Phi_W
 \le
 s_+\frac{d(d+1)}2+s_-\frac{d(d-1)}2.               \label{eq:mixed-base-cap}
\end{equation}
In particular:
\begin{enumerate}[label=(\alph*)]
\item if $W$ is transpose-stable, then $s_++s_-=w$ and
\begin{equation}
 \rank\Phi_W\le w\frac{d(d+1)}2<wd^2
 \qquad(d>1);                                      \label{eq:transpose-stable-W-cap}
\end{equation}
\item attaining the raw ordered-label rank $wd^2$ requires
\begin{equation}
 s_+=s_-=w.                                         \label{eq:both-projections-full}
\end{equation}
Hence both symmetrization and skew-symmetrization must be injective on $W$.
Equality also forces
\begin{equation}
 \dim\mathcal I_+=w\frac{d(d+1)}2,
 \qquad
 \dim\mathcal I_-=w\frac{d(d-1)}2,
 \qquad
 \mathcal I_+\cap\mathcal I_-=\{0\},             \label{eq:full-disjoint-channels}
\end{equation}
and necessarily
\begin{equation}
 W\cap W^\transpose=\{0\}.                         \label{eq:maximally-nonselfadjoint}
\end{equation}
\end{enumerate}
Thus a full-rank nonsymmetric repair cannot be a small perturbation of a
symmetric or transpose-stable family: it must be maximally
non-self-adjoint and must expose full, mutually disjoint Hermitian and
anti-Hermitian channels.  Moreover,
$W\cap W^\transpose=0$ forces
\begin{equation}
  \dim(W+W^\transpose)=2w.                        \label{eq:transpose-closure-doubles}
\end{equation}
Hence the escape doubles the dimension of the transpose-stable closure: its
canonical decomposition contains $w$ independent Hermitian and $w$
independent anti-Hermitian base directions.
\end{theorem}

\begin{proof}
Because $2$ is invertible,
$P=(P+P^\transpose)/2+(P-P^\transpose)/2$.  Apply
\Cref{thm:alignment-hermitian-factorization} to the two summands and then
sum their image-dimension bounds.  This proves
\eqref{eq:mixed-base-cap}.

If $W$ is transpose-stable, it is the direct sum of its symmetric and
skew-symmetric eigenspaces, so $s_++s_-=w$.  Since
$d(d+1)/2\ge d(d-1)/2$, \eqref{eq:transpose-stable-W-cap} follows.
For arbitrary $W$, both $s_+$ and $s_-$ are at most $w$.  Equality in the
upper bound $wd^2=w(d(d+1)/2+d(d-1)/2)$ therefore requires
\eqref{eq:both-projections-full}.  Moreover,
$\operatorname{im}\Phi_W\subseteq\mathcal I_++\mathcal I_-$, while
\[
 \dim\mathcal I_+\le s_+\frac{d(d+1)}2,
 \qquad
 \dim\mathcal I_-\le s_-\frac{d(d-1)}2.
\]
Equality $\rank\Phi_W=wd^2$ forces equality in both image bounds and in
the dimension formula for $\mathcal I_++\mathcal I_-$, proving
\eqref{eq:full-disjoint-channels}.

Suppose now that $0\ne P\in W\cap W^\transpose$.  Write $P=Q^\transpose$
with $Q\in W$.  If $P+Q\ne0$, then $P+Q$ is a nonzero symmetric element of
$W$, contradicting injectivity of skew-symmetrization.  If $P+Q=0$, then
$P$ is a nonzero skew-symmetric element of $W$, contradicting injectivity of
symmetrization.  Hence \eqref{eq:maximally-nonselfadjoint} is necessary.  Finally,
\[
 \dim(W+W^\transpose)
   =\dim W+\dim W^\transpose-\dim(W\cap W^\transpose)=2w,
\]
which proves \eqref{eq:transpose-closure-doubles}.
\end{proof}

\begin{corollary}[Alignment--Hermitian no-go theorem]
\label{cor:alignment-no-go}
Consider an odd-characteristic SNOVA-style redesign that retains:
\begin{enumerate}[label=(\roman*)]
\item degree-$d$ aligned left/right multiplier fields and a secret block
      support that is either a full $d$-dimensional linear space or has more
      than $d^2(q^{d-1}-1)$ nonzero coefficients;
\item key-randomness alignment through the standard congruence-reuse
      identity, either by pointwise commutation or by multiplier-field
      automorphism relabeling; and
\item a security count of $wd^2$ ordered power-pair features from $w$ base
      forms.
\end{enumerate}
Then no reducible or noncommutative choice of secret support, generator
change, basis change, field-size increase, distinct aligned multiplier pair,
or transpose-stable base-family modification can validate that count.  A repair
must choose at least one of the following departures:
\begin{enumerate}[label=\textbf{\arabic*.}]
\item retain symmetric or transpose-stable forms and accept the Hermitian
      cap in \eqref{eq:transpose-stable-W-cap};
\item replace them by a maximally non-self-adjoint family satisfying
      \eqref{eq:both-projections-full}--\eqref{eq:maximally-nonselfadjoint},
      including the full disjoint-channel condition
      \eqref{eq:full-disjoint-channels}, thereby restoring a full
      anti-Hermitian (exterior when $\sigma=1$) channel; or
\item abandon full alignment or change the public-key derivation so that the
      congruence-transformed forms can no longer be reused in the same way.
\end{enumerate}
The companion cryptanalysis gives explicit forgery systems for all nine
Version~2.3 rows~\cite{ThalerEtAl2026SymmetryQuotient}.  The theorem does not
assert that every construction in the second or third branch is already
broken; it proves that either branch is a fundamental architectural change
rather than a local repair of the present philosophy.
\end{corollary}

\begin{proof}
Full alignment determines the adjoint multiplier pair by
\Cref{thm:alignment-multiplier-pair}.  The symmetric and
transpose-stable branches are bounded by
\Cref{thm:alignment-hermitian-factorization,thm:mixed-base-rank-bound}.
Achieving the full ordered rank forces the second branch by
\Cref{thm:mixed-base-rank-bound}.  Refusing the adjoint-pair structure
forces abandonment of full alignment by
\Cref{thm:alignment-multiplier-pair}.
\end{proof}

The anti-Hermitian channel is not an invented bookkeeping device.  When
$\sigma=1$ it is exactly $\bigwedge^2 A$, the exterior channel that appears
in wedge and exterior-algebra cryptanalysis; the complementary Hermitian
channel is the symmetric-algebra side studied in odd-characteristic UOV
analysis~\cite{Ran2026,JinEtAl2026,BrosEtAl2026SNOVAWedge}.  Those prior
results do not by themselves prove every maximally non-self-adjoint redesign
insecure.  They do explain why restoring the missing channel is not the same
as restoring generic random-MQ behavior.

\section{Application to SNOVA Version~2.3}\label{sec:snova-application}

For one base form, the ordered label space is $A\otimes_{\F_q}A$.  The
following corollary specializes the Hermitian quotient to Version~2.3.

\begin{corollary}[Symmetric-square factorization]\label{thm:symmetry-quotient}
Let $A=\F_q[S]$ have dimension $d$, and assume $S^\transpose=S$ and
$P_i^\transpose=P_i$.  For each $i$, the map
\[
  A\otimes_{\F_q}A\longrightarrow \F_q[u]^{\rm hom}_2,
  \qquad
  \xi\otimes\eta\longmapsto u^\transpose\Sigma_\xi P_i\Sigma_\eta u
\]
is invariant under $\xi\otimes\eta\mapsto\eta\otimes\xi$.  It factors
through
\[
  \Sym^2_{\F_q}(A)=
  (A\otimes_{\F_q}A)/\Span\{\xi\otimes\eta-\eta\otimes\xi\}.
\]
Consequently the direct sum over $m_1$ base forms has rank at most
$m_1\binom{d+1}{2}$.  After any public linear map to $M$ outputs, its rank is
at most
\begin{equation}
  \min\left\{M,\;m_1\binom{d+1}{2}\right\}.          \label{eq:symmetry-cap}
\end{equation}
If the minimal polynomial of $S$ has degree $\ell$, then $d=\ell$ and the
ordered feature vector is obtained by duplicating the off-diagonal
coordinates of an unordered feature vector.
\end{corollary}

\begin{proof}
Here $\sigma=\operatorname{id}$, so the claim follows directly from
\Cref{thm:alignment-hermitian-factorization}.  Concretely, every element of
$A$ is symmetric.  Since the displayed value is a scalar,
\[
 u^\transpose\Sigma_\xi P_i\Sigma_\eta u
 =\bigl(u^\transpose\Sigma_\xi P_i\Sigma_\eta u\bigr)^\transpose
 =u^\transpose\Sigma_\eta P_i\Sigma_\xi u.
\]
Thus every alternating label $\xi\otimes\eta-\eta\otimes\xi$ lies in the
kernel.  The symmetric square has dimension $\binom{d+1}{2}$, and a later
linear map cannot enlarge its image.  When powers of $S$ form a basis, the
coordinate statement is exactly the duplication of each off-diagonal
unordered pair into its two ordered positions.
\end{proof}

The corollary gives an upper bound: extra public dependencies can only reduce
the quadratic image further.  Every concrete attack checks the exact public
rank before solving.  The outer map cannot repair the loss because it sees only
the already-equal polynomial values.  The factorization itself uses only
transposition symmetry.  Odd characteristic is needed later for the familiar
decomposition into symmetric and alternating tensors and for the polar-form
analysis.  The concrete symmetric-square corollary and the
odd-characteristic forgery analysis do not cover the submitted nonsymmetrized
$q=16$ map.  The alignment theorem still diagnoses its field-algebra
architecture, but the characteristic-two label decomposition and its wedge
attack are different.

The published construction uses a single algebra
$A=\F_q[S]\subseteq\Mat_d(\F_q)$ with $S^\transpose=S$ and symmetric base
matrices.  Full same-algebra alignment therefore has trivial adjoint on $A$,
and the Hermitian quotient is the ordinary symmetric square.  One base form
contributes at most
\[
  \binom{d+1}{2}
\]
independent quadratic features, rather than the $d^2$ ordered labels used in
the unquotiented count.  Thus
\[
 d=4:\quad 16\longrightarrow10,
 \qquad
 d=2:\quad 4\longrightarrow3.
\]
The loss occurs before any outer linear mixing and cannot be repaired by that
mixing.

The companion cryptanalysis~\cite{ThalerEtAl2026SymmetryQuotient} turns this
identity into complete square forgery systems for every published $q=19$
Version~2.3 parameter set.  That paper contains the exact affine reduction,
rejection handling, root-spectrum analysis, symbolic-homotopy recovery,
streamed-XL alternatives, and concrete numerical ledgers.  None of those
attack details are needed for the architectural theorem here; conversely, the
concrete break does not depend on the broader no-local-repair theorem.

The separation is useful.  The concrete paper answers whether the published
scheme survives.  This paper answers which parts of the surrounding design
philosophy can be retained in a repair.  Keeping full alignment and a
transpose-stable base family forces the Hermitian cap.  Recovering the raw
ordered-label count requires a maximally non-self-adjoint family and a full
anti-Hermitian channel, or else a different key-derivation mechanism.

\section{Relation to prior work and limitations}\label{sec:prior-scope}
The anti-Hermitian channel is the exterior square when the field adjoint is
trivial.  Exterior and wedge channels already play a central role in
cryptanalysis of SNOVA and related UOV-style systems
\cite{BrosEtAl2026SNOVAWedge,TsengEtAl2026Wedge}.  Complementary symmetric
channels arise in odd-characteristic UOV analyses
\cite{Ran2026,JinEtAl2026}.  Our contribution is not a new solver for those
channels.  It is a rigidity theorem explaining why full-capacity alignment
canonically produces the adjoint pair and why transpose-stable public forms
cannot support the raw ordered-label count.

Several boundaries are important.  First, a maximally non-self-adjoint family
may evade the Hermitian cap; our theorem gives necessary conditions for full
capacity, not a cryptanalysis of every family satisfying them.  Second, a
design may abandon full alignment, accept a larger public key, or use a
key-derivation identity outside the congruence-reuse model.  Third, in
characteristic two the symmetric/skew eigenspace decomposition changes and
must be analyzed through the appropriate exterior-channel formalism.  These
are genuine departures, not contradictions of the theorem.

\section{Conclusion}
Full-capacity key-randomness alignment is far more rigid than a choice of
basis or generator.  It forces degree-$d$ multiplier fields into an adjoint
pair, survives Frobenius-semilinear relabeling and large nonlinear supports,
and is uniquely maximized by the field case among irreducible secret
algebras.  In odd characteristic, the adjoint pair splits ordered quadratic
labels into Hermitian and anti-Hermitian sectors.  Symmetric or
transpose-stable public forms cannot realize the raw $d^2$ feature count.

Accordingly, a SNOVA-style repair cannot simultaneously preserve full
alignment, congruence reuse, transpose-stable forms, and the ordered-label
security rationale.  It must accept the smaller quotient, restore a full and
disjoint anti-Hermitian channel through a maximally non-self-adjoint base
family, or replace the alignment mechanism.  The companion paper shows that
the published Version~2.3 design, which takes the first branch, admits the
explicit all-nine forgery systems analyzed there.

\printbibliography
\end{document}
